% The FujiNet Palm Guide. Build with `make` (XeLaTeX); see Makefile.
\documentclass[11pt,letterpaper,openany]{book}

\usepackage[top=0.9in,bottom=0.9in,left=0.85in,right=0.85in,headheight=15pt]{geometry}
\usepackage{fontspec}
\usepackage{xcolor}
\usepackage{graphicx}
\usepackage{tikz}
\usetikzlibrary{calc,shadows,shadings,positioning,decorations.pathmorphing}
\usepackage[most]{tcolorbox}
\usepackage{titlesec}
\usepackage{fancyhdr}
\usepackage{enumitem}
\usepackage{booktabs}
\usepackage{tabularx}
\usepackage{array}
\usepackage{multicol}
\usepackage{fontawesome5}
\usepackage{microtype}
\usepackage[hidelinks]{hyperref}

% ---- Fonts ----------------------------------------------------------------
\setmainfont{Charter}
\setsansfont{Avenir Next}[BoldFont={Avenir Next Demi Bold}]
\setmonofont{Menlo}[Scale=0.85]
\newfontfamily\heavyfont{Avenir Next Heavy}
\newfontfamily\demifont{Avenir Next Demi Bold}
\newfontfamily\condfont{Avenir Next Condensed Demi Bold}
\newfontfamily\futurabold{Futura Bold}
\newfontfamily\futuramed{Futura Medium}
\newfontfamily\futuracond{Futura Condensed Medium}

% ---- Colours --------------------------------------------------------------
\definecolor{navy}{HTML}{14213D}
\definecolor{ink}{HTML}{1F2A44}
\definecolor{fujired}{HTML}{E63946}
\definecolor{nioteal}{HTML}{1A9E8F}
\definecolor{classicamber}{HTML}{E09F3E}
\definecolor{lcd}{HTML}{B9C5A6}
\definecolor{lcddark}{HTML}{8E9C7B}
\definecolor{bezel}{HTML}{3A3F47}
\definecolor{bezellight}{HTML}{6B717B}
\definecolor{paper}{HTML}{F6F4EF}
\definecolor{rule}{HTML}{C9CED6}
% The cover's 1970s palette.
\definecolor{cream}{HTML}{F2E9D8}
\definecolor{umber}{HTML}{3E2417}
\definecolor{band1}{HTML}{2C6E68}
\definecolor{band2}{HTML}{E3A72F}
\definecolor{band3}{HTML}{DC7230}
\definecolor{band4}{HTML}{B9412C}
\definecolor{band5}{HTML}{6A2C1C}

% ---- Headings, headers ------------------------------------------------------
\titleformat{\chapter}[display]
  {\sffamily\color{ink}}
  {\demifont\large\color{fujired}\MakeUppercase{\chaptertitlename}~\thechapter}
  {4pt}
  {\heavyfont\Huge}
  [\vspace{2pt}{\color{fujired}\rule{0.18\textwidth}{3pt}}]
\titlespacing*{\chapter}{0pt}{-10pt}{22pt}
\titleformat{\section}{\sffamily\bfseries\Large\color{ink}}{\thesection}{0.6em}{}
\titleformat{\subsection}{\sffamily\bfseries\large\color{ink}}{\thesubsection}{0.6em}{}
\titlespacing*{\section}{0pt}{16pt}{6pt}
\titlespacing*{\subsection}{0pt}{12pt}{4pt}
\setcounter{secnumdepth}{1}

\pagestyle{fancy}
\fancyhf{}
\renewcommand{\headrulewidth}{0pt}
\fancyhead[L]{\sffamily\small\color{bezellight}FujiNet Palm Guide}
\fancyhead[R]{\sffamily\small\color{bezellight}\leftmark}
\fancyfoot[C]{\sffamily\small\color{ink}\thepage}
\renewcommand{\chaptermark}[1]{\markboth{#1}{}}
\fancypagestyle{plain}{\fancyhf{}\fancyfoot[C]{\sffamily\small\color{ink}\thepage}}

\setlist{itemsep=3pt,topsep=4pt,parsep=0pt,partopsep=0pt}
% A PDF to read on screen: no blank pages to start chapters on the right.
\let\cleardoublepage\clearpage
\raggedbottom
\setlength{\parindent}{0pt}
\setlength{\parskip}{6pt}

% ---- Pieces ---------------------------------------------------------------
% A Palm screen, as the emulator saw it, in a small LCD bezel. #1 file, #2 caption.
\newcommand{\palmshot}[2]{%
  \begin{minipage}[t]{4.35cm}\centering
  \begin{tikzpicture}
    \IfFileExists{build/screens/#1.png}{%
      \node[inner sep=0] (s) {\includegraphics[width=3.6cm]{build/screens/#1.png}};}{%
      \node[inner sep=0,minimum size=3.6cm,fill=lcd] (s) {\sffamily\scriptsize #1};}
    \draw[bezel,line width=5pt,rounded corners=4pt]
      ($(s.south west)+(-3pt,-3pt)$) rectangle ($(s.north east)+(3pt,3pt)$);
  \end{tikzpicture}\par\vspace{3pt}
  {\sffamily\footnotesize\color{ink}#2\par}
  \end{minipage}}

% Family badges.
\newcommand{\fnbadge}{\tikz[baseline=(b.base)]\node[fill=classicamber,text=white,rounded corners=2pt,inner xsep=4pt,inner ysep=1.5pt,font=\sffamily\bfseries\footnotesize](b){FN};}
\newcommand{\fibadge}{\tikz[baseline=(b.base)]\node[fill=nioteal,text=white,rounded corners=2pt,inner xsep=4pt,inner ysep=1.5pt,font=\sffamily\bfseries\footnotesize](b){FI};}

% The title band at the top of an app page. #1 name, #2 one line, #3 badges and creators.
\newcommand{\appbanner}[3]{%
  \begin{tcolorbox}[enhanced,colback=navy,colframe=navy,arc=4pt,boxrule=0pt,
    left=10pt,right=10pt,top=8pt,bottom=8pt,
    overlay={\foreach \x in {0.55,0.63,...,0.95}{\draw[white,opacity=0.07,line width=0.6pt]
      ($(frame.north west)!\x!(frame.north east)$) -- ($(frame.south west)!\x!(frame.south east)$);}
      \fill[white] ($(frame.east)+(-1.3,0.15)$) circle (5pt);
      \fill[fujired] ($(frame.east)+(-0.9,0.15)$) circle (5pt);
      \fill[white,opacity=.35] ($(frame.east)+(-1.1,-0.2)$) circle (5pt);}]
  {\heavyfont\color{white}\LARGE #1}\par\vspace{2pt}
  {\sffamily\color{white!85!navy}#2}\par\vspace{4pt}
  {\sffamily\small\color{white}#3}
  \end{tcolorbox}}

% Side notes.
\newtcolorbox{tip}[1][Tip]{enhanced,colback=nioteal!7,colframe=nioteal,boxrule=0pt,
  leftrule=3pt,arc=0pt,left=8pt,right=6pt,top=4pt,bottom=4pt,fonttitle=\sffamily\bfseries,
  coltitle=nioteal!80!black,colbacktitle=nioteal!7,title={\faLightbulb[regular]\ #1},
  attach title to upper={\ }}
\newtcolorbox{warn}[1][Watch out]{enhanced,colback=fujired!6,colframe=fujired,boxrule=0pt,
  leftrule=3pt,arc=0pt,left=8pt,right=6pt,top=4pt,bottom=4pt,fonttitle=\sffamily\bfseries,
  coltitle=fujired!80!black,colbacktitle=fujired!6,title={\faExclamationTriangle\ #1},
  attach title to upper={\ }}

% Code and config listings.
\newtcolorbox{codebox}{colback=paper,colframe=rule,boxrule=0.5pt,arc=2pt,
  top=3pt,bottom=3pt,before upper={\setlength{\topsep}{0pt}\setlength{\partopsep}{0pt}\setlength{\parskip}{0pt}}}
\newcolumntype{L}{>{\raggedright\arraybackslash}X}

\newcommand{\btn}[1]{\tikz[baseline=(b.base)]\node[draw=ink,rounded corners=3pt,inner xsep=3pt,inner ysep=1pt,font=\sffamily\footnotesize](b){#1};}
\newcommand{\menu}[1]{\textsf{\textbf{#1}}}
\newcommand{\key}[1]{\texttt{#1}}

\begin{document}

% =========================================================================
% Cover
% =========================================================================
\begin{titlepage}
\newgeometry{margin=0pt}
\thispagestyle{empty}
\begin{tikzpicture}[remember picture,overlay,
    % The wave the serial line grows into once the byte is sent: t cm past the stop bit.
    declare function={wave(\t)=0.045*\t*\t+1.7*(\t/12.5)*sin(deg(0.72*\t));}]
  \begin{scope}[shift={(current page.south west)}]
  \fill[cream] (0,0) rectangle (21.59,27.94);

  % Masthead: series on the left, manual number on the right, a double rule.
  \node[anchor=base west,text=umber,font=\futurabold\large] at (1.6,26.0)
    {\addfontfeature{LetterSpace=18}FUJINET};
  \node[anchor=base west,text=umber!80,font=\futuracond\large] at (4.55,26.0)
    {\addfontfeature{LetterSpace=8}PALM OS SERIES};
  \node[anchor=base east,text=umber!80,font=\futuracond\large] at (20.0,26.0)
    {\addfontfeature{LetterSpace=8}MANUAL FN-PALM-1};
  \fill[umber] (1.6,25.55) rectangle (20.0,25.68);
  \fill[umber] (1.6,25.38) rectangle (20.0,25.43);

  % Title.
  \node[anchor=north west,text=umber,font=\futurabold\fontsize{60}{62}\selectfont,align=left]
    at (1.45,24.7) {The FujiNet\\Palm Guide};
  \node[anchor=north west,text=umber!85,font=\futuramed\Large,align=left,text width=17cm]
    at (1.6,19.6)
    {HotSync, TNFS and network apps for the Palm in your cradle,\\
     on FujiNet Classic and FujiNet NIO};

  % Five bands. They leave the left edge as an RS-232 frame -- start bit, the
  % letter P (0x50, least significant bit first), stop bit -- and grow into a wave.
  \def\H{9.1}\def\A{1.1}\def\bw{0.72}
  \def\steps{1.2/1,1.2/0,5.2/0,5.2/1,6.0/1,6.0/0,6.8/0,6.8/1,7.6/1,7.6/0,8.4/0,8.4/1,9.2/1}
  \def\backsteps{9.2/1,8.4/1,8.4/0,7.6/0,7.6/1,6.8/1,6.8/0,6.0/0,6.0/1,5.2/1,5.2/0,1.2/0,1.2/1,0/1}
  \foreach \c [count=\i from 0] in {band1,band2,band3,band4,band5}{
    \pgfmathsetmacro\top{\H-\i*\bw}
    \pgfmathsetmacro\bot{\H-(\i+1)*\bw}
    \fill[\c] (0,\top)
      \foreach \x/\l in \steps {-- (\x,\top-\A+\l*\A)}
      -- plot[domain=0:12.5,samples=150] ({9.2+\x},{\top+wave(\x)})
      -- plot[domain=12.5:0,samples=150] ({9.2+\x},{\bot+wave(\x)})
      \foreach \x/\l in \backsteps {-- (\x,\bot-\A+\l*\A)}
      -- cycle;
  }
  % The bits, labelled over their cells like a timing diagram.
  \foreach \x/\b in {0.6/{},1.6/ST,2.4/0,3.2/0,4.0/0,4.8/0,5.6/1,6.4/0,7.2/1,8.0/0,8.8/SP}{
    \node[text=umber!70,font=\ttfamily\scriptsize] at (\x,\H+0.38) {\b};}
  \foreach \x in {1.2,2.0,...,9.3}{
    \draw[umber!45,line width=0.4pt] (\x,\H+0.15) -- (\x,\H+0.6);}
  \node[anchor=south west,text=umber!70,font=\ttfamily\scriptsize] at (1.2,\H+0.75)
    {`P' 0x50 \textbullet\ 115200 8N1};

  % Author, date, and the FujiNet mark.
  \fill[umber] (1.6,3.55) rectangle (20.0,3.6);
  \node[anchor=north west,text=umber,font=\futurabold\Large] at (1.6,3.2) {Andrew Diller};
  \node[anchor=north west,text=umber!80,font=\futuramed\large] at (1.6,2.35) {October 2026};
  \node[anchor=north east,inner sep=0] at (20.0,3.05)
    {\includegraphics[height=1.15cm]{art/fujinet-logo.pdf}};
  \node[anchor=north east,text=umber!80,font=\futuramed\small] at (20.0,1.75) {fujinet.online};
  \end{scope}
\end{tikzpicture}
\restoregeometry
\end{titlepage}

% =========================================================================
\frontmatter
\pagestyle{plain}
\tableofcontents
\mainmatter
\pagestyle{fancy}

% =========================================================================
\chapter{Welcome}

A Palm Pilot already has everything it needs to be a network computer except
the network. It has a serial port, a cradle, and HotSync, the one-button
ritual that every Palm owner knows. \textbf{FujiNet} supplies the rest. Put a
Palm in its serial cradle, run the cradle's cable to a FujiNet, and the same
wire now does two jobs:

\begin{itemize}
  \item \textbf{HotSync.} Press the cradle's button and the Palm syncs with
    FujiNet as it once did with Palm Desktop. FujiNet installs whatever apps
    and databases are waiting in its install queue, and backs up the Palm's
    databases to its own storage.
  \item \textbf{FujiBus.} Apps on the Palm talk to FujiNet directly, at 115200
    baud, and through it to the internet: game servers, web pages, TNFS file
    servers, and friends on other Palms.
\end{itemize}

FujiNet notices which of the two is happening by itself. While an app is
talking, it leaves the line alone; when the app goes quiet, it listens now and
then for a Palm's HotSync call.

\section{Two FujiNets, two families of apps}

FujiNet comes in two firmware generations, and this guide covers both:

\begin{description}[font=\sffamily\bfseries,leftmargin=0pt]
  \item[FujiNet Classic] is the long-running \texttt{fujinet-firmware} that
    serves Atari, Apple~II, CoCo, Adam and many other machines. Its Palm
    support lives on the \texttt{palm-demo-fozz} branch. Apps written for it are
    named \textbf{FN} something, like \textsf{FN Test}. \fnbadge
  \item[FujiNet NIO] is the new firmware, \texttt{fujinet-nio}, with devices
    organised as services (network, files, WiFi, HotSync) and a portable C
    client library, \texttt{fujinet-nio-lib}. Apps written for it are named
    \textbf{FI} something, like \textsf{FI Test}. \fibadge
\end{description}

The two families do the same jobs in the same way on the Palm, and look much
alike; the difference is the FujiNet they expect at the other end of the
cable. Each FI app has its own creator code, so a Palm can carry both
families side by side, and the launcher tells them apart by name.

\begin{center}
\palmshot{fn-launcher}{The FN apps, for FujiNet Classic}\hspace{1.2cm}
\palmshot{fi-launcher}{The FI apps, for FujiNet NIO}
\end{center}

\section{The apps at a glance}

\begin{center}
\renewcommand{\arraystretch}{1.25}
\begin{tabularx}{\textwidth}{@{}l l l L@{}}
\toprule
\sffamily\bfseries Classic & \sffamily\bfseries NIO & \sffamily\bfseries Creators & \sffamily\bfseries What it does \\
\midrule
FN HotSync & FI Config & \texttt{ADHS} / \texttt{ANHS} & Browse TNFS servers and queue Palm apps for the next HotSync \\
FN Test & FI Test & \texttt{ADTS} / \texttt{ANTS} & Check the link to FujiNet, and fetch a web page \\
FN Bounce World & FI Bounce & \texttt{ADBW} / \texttt{ANBW} & A shared bouncing-shapes world on a bounce-world server \\
FN Lobby & FI Lobby & \texttt{ADLB} / \texttt{ANLB} & Every game server in the FujiNet Game Lobby \\
FN Texas Hold'em & FI Texas Hold'em & \texttt{ADTH} / \texttt{ANTH} & Multiplayer poker with Atari, Apple and CoCo players \\
FN FujiPet & FI FujiPet & \texttt{ADFP} / \texttt{ANFP} & A virtual pet that meets other Palms' pets online \\
\bottomrule
\end{tabularx}
\end{center}

All twelve run on \textbf{Palm OS 3.1 and later}: the Palm III, IIIx, IIIxe,
IIIc, V, Vx, VII, VIIx, m100 and their cousins. Chapter~\ref{ch:compat} has the
details.

% =========================================================================
\chapter{Classic and NIO}
\label{ch:diff}

Both firmwares give the Palm the same two services on one cable, HotSync and
FujiBus, and both share the line the same way. They differ in how FujiNet is
configured, where it keeps the install queue, and what the apps speak on the
wire.

\section{Side by side}

{\small
\renewcommand{\arraystretch}{1.3}
\begin{tabularx}{\textwidth}{@{}>{\raggedright\sffamily\bfseries}p{2.6cm} L L@{}}
\toprule
 & \fnbadge\ \sffamily\bfseries Classic & \fibadge\ \sffamily\bfseries NIO \\
\midrule
Firmware & \texttt{fujinet-firmware}, branch \texttt{palm-demo-fozz}; board
  \texttt{fujinet-rs232-s3-palm} (FN-RS232) or FujiNet-PC
  & \texttt{fujinet-nio} with HotSync (pull request \#19); the FN-RS232 ESP32-S3
  build, or the POSIX build on a computer \\
Apps & FN \emph{name}, creator \texttt{AD..} & FI \emph{name}, creator \texttt{AN..} \\
Talking to FujiNet & FujiBus commands from the app's own small client
  (\texttt{fnlink}); network through the \texttt{N1:} device with devicespecs
  like \texttt{N1:TCP://host:port/}
  & \texttt{fujinet-nio-lib}: FujiDevice, NetworkDevice, FileDevice and WiFi
  services; a lost reply is asked for again and NIO replays it \\
HotSync settings & \texttt{[HotSync]} in \texttt{fnconfig.ini}, or the web UI's
  \textsf{Palm HotSync} panel & \texttt{hotsync:} in \texttt{fujinet.yaml}, or
  \texttt{hotsync.set} and \texttt{hotsync.save} on the console \\
Install queue & \texttt{/palm/install/} on the SD card & \texttt{hotsync/install/}
  on the persistent filesystem (the SD card on an FN-RS232) \\
Backups & \texttt{/palm/backup/\emph{user}/} & \texttt{hotsync/backup/\emph{user}/} \\
Queueing an app & From the Palm: \textsf{FN HotSync} copies a file from a TNFS
  host slot. From a computer: WebDAV, \texttt{curl -T app.prc
  http://fujinet/dav/palm/install/}
  & From the Palm: \textsf{FI Config} copies a file from a TNFS server. From the
  console: \texttt{hotsync.fetch http://host/app.prc} \\
Network HotSync & TCP 14238, and 6416 for emulators & TCP 14238 (NetSync) \\
Calendar & Copies a Google or iCalendar calendar into the Date Book at each sync
  & Not yet \\
Apps on TNFS & \texttt{fujinet.diller.org/PALM/} & \texttt{fujinet.diller.org/NIO/PALM/} \\
\bottomrule
\end{tabularx}}

\section{Sharing one cable}

A HotSync begins at 9600 baud, while apps talk at 115200. FujiNet cannot hear
a 9600-baud call while it is listening at 115200, so both firmwares do the
same dance:

\begin{enumerate}
  \item When no app has spoken for a minute, FujiNet starts \emph{listening
    windows}: it drops the line to 9600 baud for 1.5~seconds, listens for a
    Palm's HotSync wake-up, then goes back to 115200.
  \item Any app request holds the windows off for another minute, so a running
    app is never interrupted.
  \item Bytes that are not an app request end the wait at once: that is what a
    pressed HotSync button sounds like at 115200.
  \item Once the Palm's wake-up is answered, the sync owns the line until it
    ends, including the switch to a faster rate.
\end{enumerate}

An app request that lands inside a window gets no answer, so every app
resends. On NIO, \texttt{fujinet-nio-lib} also wakes the line with a harmless
request of its own after an idle spell. You will not notice any of this,
except that the first request after a long pause can take a second longer.

\begin{tip}
A Palm running an app cannot HotSync, and a Palm that is HotSyncing is not
running an app. Leave the app (or use its \btn{HotSync} button) before you
press the cradle button.
\end{tip}

\section{Which one do I have?}

Run \textsf{FN Test} and \textsf{FI Test} and tap \btn{Ping}. The one that
says \textsf{FujiNet found!} matches your FujiNet. \textsf{FI Test} also
reports the board NIO is running on, such as \textsf{POSIX + FujiBus over
RS-232} for a computer running NIO.

% =========================================================================
\chapter{Getting set up}
\label{ch:setup}

\section{Hardware}

\begin{itemize}
  \item A Palm OS handheld (OS 3.1 or later) and its \textbf{serial} HotSync
    cradle or HotSync cable, the kind with a DB-9 plug. USB cradles do not
    work.
  \item A FujiNet with an RS-232 port: the \textbf{FN-RS232} board, or a
    computer running FujiNet-PC or FujiNet NIO with a serial port.
  \item A \textbf{null-modem adapter} between the cradle's DB-9 and the
    FN-RS232's DB-9. Both ends are wired as a computer would be.
\end{itemize}

\begin{warn}
A Palm cradle never raises the RTS/CTS flow-control lines. Build or configure
FujiNet without hardware flow control, or it will wait forever for permission
to send. The classic \texttt{fujinet-rs232-s3-palm} board and NIO's default
\texttt{flow\_control: none} already do.
\end{warn}

\section{FujiNet Classic}

Flash the \texttt{fujinet-rs232-s3-palm} build of the \texttt{palm-demo-fozz}
branch. HotSync is on by default; its settings are in \texttt{fnconfig.ini}
and in the web UI's \textsf{Palm HotSync} panel:

\begin{codebox}
\begin{verbatim}
[HotSync]
enabled=1
user=FujiNet          ; given to a Palm that has never synced
backup=flagged        ; none | flagged | all
netsync_port=14238    ; 0 turns it off
serial_port=bus       ; the cradle shares the RS-232 bus line
\end{verbatim}
\end{codebox}

An SD card is required: the queue is \texttt{/palm/install/}, installed files
move to \texttt{/palm/installed/}, and backups go to
\texttt{/palm/backup/\emph{user}/}. Classic can also fill the Palm's Date Book
from a Google or iCalendar calendar at each sync (\texttt{calendar=GCAL:///}
or \texttt{ICAL://host/feed.ics}).

\section{FujiNet NIO}

HotSync is off by default. Turn it on in \texttt{fujinet.yaml}, or from the
diagnostics console, and restart:

\begin{codebox}
\begin{verbatim}
hotsync:
  enabled: true
  cradle: true         # a Palm cradle shares the FujiBus serial line
  port: 14238          # NetSync TCP port; 0 turns it off
  root: /hotsync       # folder on the persistent filesystem
  user_name: FujiNet   # given to a Palm that has never synced
  backup: flagged      # none | flagged | all
\end{verbatim}
\end{codebox}

\begin{codebox}
\begin{verbatim}
hotsync.set enabled 1
hotsync.set cradle 1
hotsync.save
hotsync.status        # is it listening? what did the last sync do?
\end{verbatim}
\end{codebox}

\section{Installing the apps}

The first app has to arrive by HotSync, so put it in the install queue and
press the cradle button. After that, \textsf{FN HotSync} or \textsf{FI
Config} can queue the rest from the Palm itself. The builds are on the TNFS
server \texttt{fujinet.diller.org}: the FN apps in \texttt{/PALM/}, the FI apps
in \texttt{/NIO/PALM/}.

\begin{description}[font=\sffamily\bfseries,leftmargin=0pt]
  \item[Classic] Copy the \texttt{.prc} to \texttt{/palm/install/} on the SD
    card, or send it over WebDAV:\\
    \texttt{curl -T FN-HotSync.prc http://fujinet.local/dav/palm/install/}
  \item[NIO] Serve the \texttt{.prc} from any computer and fetch it from the
    console:\\
    \texttt{python3 -m http.server 8000} \quad (on the computer)\\
    \texttt{hotsync.fetch http://192.168.1.5:8000/FI-Config.prc} \quad (on NIO)
\end{description}

Then press the cradle's HotSync button. The Palm shows its HotSync screen,
FujiNet installs the queued files and moves them to \texttt{installed/}, and
the new icons appear in the launcher.

\begin{center}
\palmshot{fi-hotsync}{A cradle HotSync with NIO installing a queued app}
\end{center}

% =========================================================================
\chapter{Palm OS versions}
\label{ch:compat}

Every app runs on Palm OS 3.1 and later. Palm OS 3.3 brought a new Serial
Manager; on 3.1 and 3.2 the apps (and \texttt{fujinet-nio-lib}) fall back to
the older serial library, which does the same job. An app started on Palm OS
3.0 or earlier says \textsf{This app needs Palm OS 3.1 or later} and quits.

FN FujiPet and FI FujiPet draw in 256 colours on a colour Palm running OS 3.5
or later, and in black and white with a grey dither everywhere else. The other
apps are black and white on every Palm.

\begin{center}
\renewcommand{\arraystretch}{1.25}
\begin{tabular}{@{}l l l l@{}}
\toprule
\sffamily\bfseries Palm OS & \sffamily\bfseries Devices & \sffamily\bfseries Serial & \sffamily\bfseries Tested \\
\midrule
3.0 & Palm III & --- & the apps refuse to start \\
3.1 & Palm IIIx, Palm V & serial library & emulator: all twelve apps \\
3.2 & Palm VII & serial library & no ROM available \\
3.3 & Palm III, V (upgraded) & Serial Manager & emulator, and a real Palm in a cradle \\
3.5 & IIIxe, IIIc, Vx, VIIx, m100 & Serial Manager & emulator (V, VIIx, m100) \\
4.x & m500, m505, IIIc & Serial Manager & emulator (IIIc) \\
\bottomrule
\end{tabular}
\end{center}

\begin{tip}[Newer Palms]
Palms from the m500 series onward use a different cradle connector, and their
serial cradles are rare. Anything with Palm OS 3.1 or later and a serial
cradle works.
\end{tip}

% =========================================================================
\chapter{The apps}
\label{ch:apps}

Each app has a page of its own. Where the FN and FI versions differ, the page
says how; otherwise everything on it holds for both.

\begin{center}
\renewcommand{\arraystretch}{1.6}
\begin{tabularx}{\textwidth}{@{}l l X r@{}}
\toprule
\sffamily\bfseries \fnbadge\ Classic & \sffamily\bfseries \fibadge\ NIO & \sffamily\bfseries For & \sffamily\bfseries Page \\
\midrule
FN HotSync & & queueing apps from TNFS host slots & \pageref{app:fnhotsync} \\
 & FI Config & queueing apps from TNFS servers & \pageref{app:ficonfig} \\
FN Test & FI Test & checking the cradle link & \pageref{app:test} \\
FN Bounce World & FI Bounce & a shared bouncing world & \pageref{app:bounce} \\
FN Lobby & FI Lobby & the FujiNet Game Lobby & \pageref{app:lobby} \\
FN Texas Hold'em & FI Texas Hold'em & multiplayer poker & \pageref{app:holdem} \\
FN FujiPet & FI FujiPet & a virtual pet with friends & \pageref{app:fujipet} \\
\bottomrule
\end{tabularx}
\end{center}

\begin{tip}[Finding them]
New apps land in the launcher's \textsf{Unfiled} category. Pick
\textsf{Unfiled} or \textsf{All} from the category list at the top right of
the launcher to see them.
\end{tip}

% One page per app; \input by fujinet-palm-guide.tex.

% ---------------------------------------------------------------------------
\newpage
\section{FN HotSync}\label{app:fnhotsync}
\appbanner{FN HotSync}{Browse FujiNet's host slots and queue Palm apps for the next HotSync}
  {\fnbadge\quad creator \texttt{ADHS}\quad FujiNet Classic}

FN HotSync is the Palm's own window on FujiNet Classic's host slots: the
same eight slots an Atari or Apple sees in CONFIG, one of them FujiNet's SD
card and the rest TNFS servers. Browse a server, pick a Palm app or database,
and FN HotSync has FujiNet copy it into \texttt{/palm/install/} on the SD card.
The next HotSync installs it.

\begin{center}
\palmshot{fn-hotsync-main}{The host slots}\hfill
\palmshot{fn-hotsync-browse}{Browsing \texttt{/PALM/} on a TNFS server}\hfill
\palmshot{fn-hotsync-queue}{Queue a file for HotSync}
\end{center}

\begin{multicols}{2}
\subsection*{Using it}
\begin{itemize}
  \item Pick a host slot and tap \btn{Browse}. Folders end in \texttt{/}: select
    one and tap \btn{Open} to go in, \btn{Up} to come back out.
  \item Select a \texttt{.prc}, \texttt{.pdb} or \texttt{.pqa} file and tap
    \btn{Open}; confirm with \btn{Queue}. FujiNet copies it to the SD card.
  \item \btn{Queue} on the main screen lists what is waiting;
    \btn{HotSync Now} there starts the sync.
  \item \btn{Edit} names a slot: \texttt{SD} for FujiNet's card, or a TNFS
    server such as \texttt{apps.irata.online}. FujiNet saves the slots.
  \item \btn{Info} shows FujiNet's WiFi network, IP address and firmware
    version.
  \item \btn{HotSync} closes the cradle port and opens the Palm's HotSync
    app, ready to sync.
\end{itemize}

\subsection*{Good to know}
\begin{itemize}
  \item Queued files go to the slot named \texttt{SD}; without one, FN
    HotSync says so.
  \item The copy runs on FujiNet, server to SD card, so even a large app
    takes only a moment.
  \item A queued file stays until a HotSync installs it; the Palm cannot take
    it back out of the queue.
  \item The NIO version of this app is FI Config, on the next page.
\end{itemize}
\end{multicols}

\begin{center}
\palmshot{fn-hotsync-edit}{Naming a host slot}\hfill
\palmshot{fn-hotsync-queued}{What is waiting for the next sync}\hfill
\palmshot{fn-hotsync-info}{FujiNet's network}
\end{center}

% ---------------------------------------------------------------------------
\newpage
\section{FI Config}\label{app:ficonfig}
\appbanner{FI Config}{Browse TNFS servers and queue Palm apps for HotSync through FujiNet NIO}
  {\fibadge\quad creator \texttt{ANHS}\quad FujiNet NIO}

FI Config is FN HotSync rebuilt for FujiNet NIO. NIO has no host slots, so
FI Config keeps its own list of eight TNFS servers, starting with
\texttt{fujinet.diller.org} and \texttt{apps.irata.online}. Browse one, pick
an app, and FI Config copies it through NIO's file service into
\texttt{hotsync/install/}, ready for the next HotSync.

\begin{center}
\palmshot{fi-config-main}{The TNFS servers}\hfill
\palmshot{fi-config-browse}{Browsing a server}\hfill
\palmshot{fi-config-progress}{Copying into the queue}
\end{center}

\begin{multicols}{2}
\subsection*{Using it}
\begin{itemize}
  \item Pick a server and tap \btn{Browse}; \btn{Open} goes into a folder or
    picks a file, \btn{Up} goes back.
  \item Opening a \texttt{.prc}, \texttt{.pdb} or \texttt{.pqa} asks whether
    to queue it. The copy shows its progress on the line above the list.
  \item \btn{Queue} lists the files waiting in NIO's install folder.
  \item \btn{Edit} changes a server's name; clear it to empty the slot. The
    list lives in the app's own preferences.
  \item \btn{Info} shows NIO's WiFi network, IP address, firmware and board,
    and how many replies were asked for again.
  \item \btn{HotSync} hands the cradle back and opens the HotSync app.
\end{itemize}

\subsection*{Good to know}
\begin{itemize}
  \item The Palm does the copying, a few hundred bytes at a time, so a 30~KB
    app takes about half a minute over TNFS.
  \item If a copy is cut short, HotSync refuses the half-written file rather
    than installing a broken app: queue it again to replace it.
  \item A reply lost on the cable is simply asked for again; Info counts how
    often that happened.
  \item On NIO's console, \texttt{hotsync.fetch} queues an app over HTTP
    without the Palm.
\end{itemize}
\end{multicols}

\begin{center}
\palmshot{fi-config-queued}{The install queue}\hfill
\palmshot{fi-config-info}{NIO's firmware and network}\hfill
\palmshot{fi-hotsync-log}{The HotSync log: installed}
\end{center}

% ---------------------------------------------------------------------------
\newpage
\section{FN Test and FI Test}\label{app:test}
\appbanner{FN Test \textbullet\ FI Test}{Is FujiNet there? A one-tap check of the cradle link}
  {\fnbadge\ \texttt{ADTS}\qquad \fibadge\ \texttt{ANTS}}

The first app to run on a new setup. \btn{Ping} sends FujiNet one request
and shows what comes back. If it says \textsf{FujiNet found!}, the cradle,
the cable, the null-modem adapter and the firmware are all doing their job,
and every other app will work too.

\begin{center}
\palmshot{fn-test}{FN Test finds FujiNet Classic}\hfill
\palmshot{fi-test-ping}{FI Test finds FujiNet NIO}\hfill
\palmshot{fi-test-fetch}{FI Test fetches a web page}
\end{center}

\begin{multicols}{2}
\subsection*{FN Test}
\btn{Ping FujiNet} asks FujiNet Classic for its adapter configuration and
shows the firmware version, FujiNet's IP address and the WiFi network it is
on.

\subsection*{FI Test}
\begin{itemize}
  \item \btn{Ping} asks NIO for its version and board (FujiDevice), then
    for its WiFi status and settings.
  \item \btn{Fetch} gets \texttt{http://example.com/} through NIO's network
    service and shows the HTTP status, the size and the page's title: proof
    that the Palm can reach the internet.
\end{itemize}

\columnbreak
\subsection*{If it fails}
\textsf{No answer from FujiNet} means nothing came back after several tries.
Check the cable and adapter, and that you are running the right app for your
FujiNet: FN Test for Classic, FI Test for NIO. Chapter~\ref{ch:diff} explains
the difference.

\begin{tip}
Right after a HotSync, or after a long pause, the first ping can take a
second longer: FujiNet may be listening for a HotSync at that moment, and the
app simply asks again.
\end{tip}
\end{multicols}

% ---------------------------------------------------------------------------
\newpage
\section{Bounce World}\label{app:bounce}
\appbanner{FN Bounce World \textbullet\ FI Bounce}{Shapes bouncing across a world shared by every connected machine}
  {\fnbadge\ \texttt{ADBW}\qquad \fibadge\ \texttt{ANBW}}

A client for Mark Fisher's \textsf{bounce-world} server. Every machine that
connects (an Atari, a C64, a Palm) gets a window onto one shared world, and
the shapes in it bounce from screen to screen. The Palm draws its part of the
world in block graphics, polling the server several times a second, which
also makes it a fine workout for the cradle link.

\begin{center}
\palmshot{fn-bounce-connect}{Connecting to a server}\hfill
\palmshot{fi-bounce-world}{The world in motion}\hfill
\palmshot{fi-bounce-info}{FI Bounce's link counters}
\end{center}

\begin{multicols}{2}
\subsection*{Using it}
Enter the \textsf{Server} and \textsf{Port} (9003 for the framed protocol of
bounce-world 2.3 and later) and a \textsf{Name} of up to 8 letters, then tap
\btn{Connect}. FI
Bounce starts with \texttt{bouncy.diller.org}; FN Bounce World with
\texttt{localhost}, for FujiNet-PC.

\begin{tabularx}{\linewidth}{@{}l X@{}}
\btn{Frz} \key{F} & freeze or unfreeze the world, for everyone \\
\btn{Rst} \key{R} & reset it \\
\btn{+} \btn{-} & faster, slower (also the scroll keys) \\
\btn{Add} \key{1}--\key{5} & add a body of that size \\
\btn{Who} \key{W} & who else is connected \\
\btn{Info} \key{I} & your name, clients, world size and bodies \\
\key{Q} & disconnect \\
\end{tabularx}

While you are in the world the app buttons play too: 1 freezes, 2 adds a
body, 3 shows who, 4 shows info.

\subsection*{FI Bounce's counters}
FI Bounce's Info adds a line of link statistics: requests, exchanges with
NIO, retries and replies lost. A long run with no losses is a healthy cable.

\begin{warn}[Known issue]
The world keeps drawing while the \menu{World} menu is open, so the menu
shows as a blank box. Its items still work, but the buttons and Graffiti keys
are easier.
\end{warn}
\end{multicols}

% ---------------------------------------------------------------------------
\newpage
\section{Lobby}\label{app:lobby}
\appbanner{FN Lobby \textbullet\ FI Lobby}{Every game server in the FujiNet Game Lobby, in your pocket}
  {\fnbadge\ \texttt{ADLB}\qquad \fibadge\ \texttt{ANLB}}

The FujiNet Game Lobby lists the multiplayer games FujiNet machines can
join: Battleship, 5 Card Stud, Texas Hold'em and many more, with who is
playing and from which platforms. The Lobby app reads the whole list,
\texttt{/viewFull} from \texttt{lobby.fujinet.online}, and shows it on the
Palm.

\begin{center}
\palmshot{fi-lobby-loading}{Reading the lobby}\hfill
\palmshot{fn-lobby-list}{Every server, by game}\hfill
\palmshot{fn-lobby-detail}{One server in detail}
\end{center}

\begin{multicols}{2}
\subsection*{Using it}
\begin{itemize}
  \item The list loads when the app starts. Each row is a game, a server
    and its players, such as \textsf{0/4}.
  \item Tap a row for the server's status, region, players, app key and URL,
    and the platforms it serves with the client URL for each. The scroll keys
    page through it.
  \item \btn{Refresh} reads the list again.
  \item \menu{Lobby > Use QA Lobby} switches to the test lobby, for
    developers; choose it again to go back.
\end{itemize}

\subsection*{Good to know}
\begin{itemize}
  \item The reply is about 30~KB of JSON, more than a Palm would like to hold
    at once. The app reads it as it streams in, keeping only what it shows,
    up to 48 servers.
  \item FI Lobby stays usable while the list loads: the menus, Refresh and
    Details all answer straight away.
\end{itemize}
\end{multicols}

% ---------------------------------------------------------------------------
\newpage
\section{Texas Hold'em}\label{app:holdem}
\appbanner{FN Texas Hold'em \textbullet\ FI Texas Hold'em}{Multiplayer poker, at the same tables as Atari, Apple II and CoCo players}
  {\fnbadge\ \texttt{ADTH}\qquad \fibadge\ \texttt{ANTH}}

A Palm client for the FujiNet Texas Hold'em server, listed in the Game Lobby.
Sit at a table with AI players or with friends on other machines; the server
deals and runs every rule, and the Palm shows the felt, the cards and the
moves you can make.

\begin{center}
\palmshot{fn-holdem-setup}{Pick a name and a table}\hfill
\palmshot{fn-holdem-table}{Your move}\hfill
\palmshot{fi-holdem-showdown}{The showdown}
\end{center}

\begin{multicols}{2}
\subsection*{Using it}
\begin{itemize}
  \item Enter a \textsf{Name} (up to 8 letters and digits). The tables on the
    \textsf{Server} are listed with their players; pick one (or type its name
    under \textsf{Table}) and tap \btn{Sit Down}.
  \item Your two cards sit beside the board; the other players are above, face
    down until the showdown.
  \item On your turn the moves appear along the bottom with the server's
    clock. Tap one, or pick one with the scroll keys and press any app button.
  \item Graffiti works too: \key{f} fold, \key{c} check or call, \key{b} bet,
    \key{r} raise, \key{a} all-in.
  \item Tap the status line, or \menu{Table > Last Result}, for the last
    hand. \menu{Table > Leave Table} gives up your seat.
\end{itemize}

\subsection*{Good to know}
\begin{itemize}
  \item The default server is \texttt{th.carr-designs.com}. A lobby address
    that ends in \texttt{?table=den} can be typed as the server, and the table
    is taken from it.
  \item The cards are the Apple~II client's art, which needs no colour; the
    small cards are the Atari Lynx client's.
  \item Leaving the app, or switching to another, leaves the table for you.
\end{itemize}
\end{multicols}

% ---------------------------------------------------------------------------
\newpage
\section{FujiPet}\label{app:fujipet}
\appbanner{FN FujiPet \textbullet\ FI FujiPet}{A virtual pet that grows up like the 1996 original, and makes friends online}
  {\fnbadge\ \texttt{ADFP}\qquad \fibadge\ \texttt{ANFP}}

Name your pet and pick a Cat, Dog, Bunny or Dino; the egg hatches five
minutes later. It grows from baby to child, teen and adult over a few days,
and how well it is looked after decides what it grows into and how long it
lives. With \textsf{Meet friends online} on, it keeps a connection to a
FujiPet friend server and can chat with pets on other Palms.

\begin{center}
\palmshot{fn-fujipet-new}{A new egg}\hfill
\palmshot{fn-fujipet-main}{All grown up, in colour}\hfill
\palmshot{fn-fujipet-friends}{Friends online}
\end{center}

\begin{multicols}{2}
\subsection*{Looking after it}
The eight icons are the original's. Along the top: \textbf{Feed} (a meal
fills a hunger heart, a snack a happy one), \textbf{Light}, \textbf{Play}
(guess which way it will look, three of five to win) and \textbf{Medicine}.
Along the bottom: \textbf{Clean}, \textbf{Status}, \textbf{Scold} and
\textbf{Friends}. Tap the pet itself to make it happy. A heart running empty makes it call
(the \textsf{!} in the corner); a call left for 15 minutes is a care mistake.
Mess and too many sweets make it sick. The pet lives on while the app is
closed, and catches up when you return.

\subsection*{Friends}
\btn{Register} adds your pet to the friend server. The Friends screen lists
the pets there, online ones first in bold; pick one, choose one of eight
messages under \textsf{Say} and tap \btn{Send}. It pops up as a bubble on
their Palm at once, and a kind one cheers their pet up.

\subsection*{Settings}
\menu{Options > Settings} sets the friend server (\textsf{host:port}), turns
going online and sounds on or off, and runs time \textsf{Fast} (a pet minute
a second) or \textsf{Turbo} (a pet day a minute) to watch it grow up.
\end{multicols}

\begin{center}
\palmshot{fn-fujipet-status}{Status: hearts, age and weight}\hfill
\palmshot{fi-fujipet-message}{A message from a friend}\hfill
\palmshot{fn-fujipet-mono}{In greys on a Palm V}
\end{center}


% =========================================================================
\chapter{When things go wrong}

\begin{description}[font=\sffamily\bfseries\color{fujired},leftmargin=0pt,style=nextline,itemsep=8pt]
  \item[\textsf{No answer from FujiNet. Is the cradle cabled to it?}]
    The app sent its request several times and heard nothing. Check the
    null-modem adapter, that FujiNet is running, that it was built without
    hardware flow control, and that you are running the right family: FN apps
    with Classic, FI apps with NIO.
  \item[\textsf{This app needs Palm OS 3.1 or later}]
    The Palm runs Palm OS 3.0 or earlier.
  \item[The HotSync button does nothing]
    Quit any FujiNet app first. Check that HotSync is on: on NIO,
    \texttt{hotsync.status} should say \texttt{enabled: 1} and
    \texttt{listening: 1}; on Classic, \texttt{[HotSync] enabled=1} with
    \texttt{serial\_port=bus}. Give it a few seconds: FujiNet listens in
    windows.
  \item[A queued app did not install]
    Look in the queue (\btn{Queue} in FN HotSync or FI Config). On NIO,
    \texttt{hotsync.status} reports \texttt{last\_installed} and
    \texttt{last\_install\_failures}. A copy that FI Config could not finish is
    refused by HotSync; queue it again.
  \item[FujiPet says \textsf{No friend server}]
    Check the server's \textsf{host:port} in \menu{Options > Settings}. A
    computer running NIO on a Mac connects to the IPv6 address for
    \texttt{localhost}, so start the friend server with \texttt{-{}-host \char34\char34}, or
    give its IPv4 address.
  \item[A FujiNet-PC on a Mac drops requests]
    FujiNet-PC builds before the \texttt{[pc]} serial fix left the terminal's
    extended input processing on, which swallows bytes \texttt{0x16} and
    \texttt{0x0F} in a packet. Use a build with the fix.
\end{description}

% =========================================================================
\appendix
\chapter{For developers}

\section{Building}

The apps live in \texttt{github.com/dillera/fujinet-palm}: the FN apps on the
\texttt{classic} branch, the FI apps on \texttt{nio}. They build with
prc-tools-remix and PilRC in a Docker image:

\begin{codebox}
\begin{verbatim}
make image    # once: the palmdev toolchain image
make          # every app, apps/*/*.prc
make test     # desktop tests of the shared C code
make traps    # Palm OS calls newer than 3.1, per app
\end{verbatim}
\end{codebox}

On \texttt{nio}, \texttt{make} also builds \texttt{fujinet-nio-lib}'s
\texttt{palm} target from \texttt{../fujinet-nio-workspace}. An FI app sets
\texttt{NIO = 1} in its Makefile; \texttt{fn\_read} never blocks, so poll it
from the event loop, and when a reply is lost ask again at the same offset.

\section{Testing in an emulator}

\texttt{assets/} holds Palm OS ROMs and CloudpilotEmu images with Setup done,
for Palm OS 3.1, 3.3, 3.5 and 4.0, and the patch that gives CloudpilotEmu a
\texttt{-{}-serial} option and scripting commands. \texttt{tools/rig.sh} joins
an emulated Palm's cradle to FujiNet-PC or to FujiNet NIO through a pair of
pseudo-terminals:

\begin{codebox}
\begin{verbatim}
cp assets/images/palm-m100-os3.5.1.img runs/t/palm.img
tools/rig.sh start runs/t runs/t/palm.img nio
tools/rig.sh cmd runs/t install $PWD/apps/FujiNetTest/FujiNetTest.prc
tools/rig.sh cmd runs/t launch 'FI\ Test'
tools/rig.sh cmd runs/t tap 40 147
tools/rig.sh cmd runs/t screenshot ping.png
\end{verbatim}
\end{codebox}

Every screen in this guide was taken that way.

\end{document}
